% =============================================================================
% Research Paper: SPIN Sales Coach — Conversational AI Agent
% MSc DSAI - M1 Internship Report (Agent Component)
% =============================================================================
\documentclass[10pt,twocolumn,a4paper]{article}

% --- Packages ---
\usepackage[utf8]{inputenc}
\usepackage[T1]{fontenc}
\usepackage{lmodern}
\usepackage[english]{babel}
\usepackage{geometry}
\geometry{top=2cm, bottom=2cm, left=1.8cm, right=1.8cm}
\usepackage{graphicx}
\usepackage{booktabs}
\usepackage{multirow}
\usepackage{float}
\usepackage{amsmath}
\usepackage{amssymb}
\usepackage{hyperref}
\usepackage{xcolor}
\usepackage{caption}
\usepackage{subcaption}
\usepackage[numbers,sort&compress]{natbib}
\usepackage{enumitem}
\usepackage{tikz}
\usepackage{fancyhdr}
\usepackage{titlesec}
\usepackage{listings}

% --- Formatting ---
\hypersetup{
    colorlinks=true,
    linkcolor=blue!60!black,
    citecolor=blue!60!black,
    urlcolor=blue!60!black
}

\setlength{\columnsep}{0.6cm}
\titlespacing*{\section}{0pt}{1.2ex plus 0.5ex minus 0.2ex}{0.8ex plus 0.2ex}
\titlespacing*{\subsection}{0pt}{1.0ex plus 0.4ex minus 0.2ex}{0.6ex plus 0.2ex}
\titleformat{\section}{\normalfont\large\bfseries}{\thesection.}{0.5em}{}
\titleformat{\subsection}{\normalfont\normalsize\bfseries}{\thesubsection}{0.5em}{}

\pagestyle{fancy}
\fancyhf{}
\fancyhead[L]{\small\textit{SPIN Sales Coach Agent --- Conversational AI Evaluation}}
\fancyhead[R]{\small\thepage}
\renewcommand{\headrulewidth}{0.4pt}

\lstset{basicstyle=\ttfamily\scriptsize, breaklines=true, frame=single}

% --- Title ---
\title{\vspace{-1.5cm}\textbf{Design, Deployment, and Quantitative Evaluation of a\\Conversational AI Sales Coach Using Google ADK\\and Large Language Models}\vspace{-0.3cm}}

\author{
    \textbf{Ezhilarasi Muthukumar}\\
    \small MSc Data Science \& Artificial Intelligence (M1)\\
    \small Internship at a veterinary pharmaceutical company, France\\
    \small \texttt{ezhilarasi.muthukumar@etu.univ-cotedazur.fr}
}
\date{\small July 2026}

\begin{document}
\maketitle
\thispagestyle{fancy}

% =============================================================================
% ABSTRACT
% =============================================================================
\begin{abstract}
\noindent This paper presents the design, cloud deployment, and quantitative evaluation of a domain-specific conversational AI agent built using Google's Agent Development Kit (ADK) and Gemini large language models. The agent --- a SPIN Selling methodology coach for pharmaceutical sales representatives --- operates in three modes (teaching, evaluation, role-play) with Retrieval-Augmented Generation (RAG) grounding from a curated knowledge base. The system is deployed on Google Cloud Run and validated through an automated evaluation framework comprising 60 curated test cases across 6 quality categories. Deterministic evaluation metrics demonstrate 77\% tool-selection accuracy, 75\% keyword coverage, and 0\% runtime errors across all test prompts. The agent consistently maintains its coaching persona, correctly identifies SPIN phases, and provides actionable domain-specific guidance. A synthetic knowledge base was used exclusively for validation; the production system maps to internal operational data via RAG. This work presents a rigorously described and empirically validated contribution to domain-specific conversational AI by describing the technical architecture, orchestration framework, and evaluation methodology.

\medskip
\noindent\textbf{Keywords:} Conversational AI, LLM Agent, Google ADK, Gemini, SPIN Selling, RAG, Sales Coaching, Cloud Run, Evaluation, LLM-as-Judge
\end{abstract}

% =============================================================================
% 1. INTRODUCTION
% =============================================================================
\section{Introduction}
\label{sec:intro}

The emergence of large language models (LLMs) has enabled a new class of task-oriented conversational agents that can operate as domain-specific experts~\cite{wei2022emergent}. Unlike general-purpose chatbots, these agents maintain persistent personas, ground responses in retrieved knowledge, and use structured tools to perform actions~\cite{schick2023toolformer}.

This work addresses a concrete business need: training pharmaceutical sales representatives in the SPIN Selling methodology~\cite{rackham1988spin} --- a structured questioning technique for consultative sales. The agent serves as an always-available coach that teaches methodology, evaluates practice attempts, and conducts role-play scenarios with simulated buyer personas.

The contribution of this work is to: (1)~rigorously describe the technical architecture, (2)~identify the LLM models and orchestration framework, (3)~define metrics for evaluating conversational quality, and (4)~produce quantitative prototype evaluation.

\subsection{Research Questions}
\begin{enumerate}[leftmargin=*,noitemsep,topsep=2pt]
    \item[\textbf{RQ1:}] Can a single-agent architecture with RAG maintain domain expertise and role consistency across diverse interaction modes?
    \item[\textbf{RQ2:}] What quantitative metrics effectively evaluate conversational quality for a coaching agent?
    \item[\textbf{RQ3:}] How does the agent perform on automated evaluation across teaching, evaluation, role-play, and adversarial scenarios?
\end{enumerate}

% =============================================================================
% 2. TECHNICAL ARCHITECTURE
% =============================================================================
\section{Technical Architecture}
\label{sec:architecture}

\subsection{System Overview}

The agent is built on Google's Agent Development Kit (ADK)~\cite{googleadk2024}, an open-source Python framework for constructing LLM-powered agents with tool use and session management.

\begin{table*}[t]
\caption{System architecture components.}
\label{tab:architecture}
\setlength{\tabcolsep}{36pt}
\renewcommand{\arraystretch}{1.4}
\begin{tabular}{l l}
\toprule
\textbf{Component} & \textbf{Technology} \\
\midrule
LLM Model & Gemini 3.1 Flash Lite \\
Orchestration Framework & Google ADK (v0.3+) \\
RAG Knowledge Base & Custom JSON (keyword retrieval) \\
Session Management & ADK InMemorySessionService \\
Deployment & Google Cloud Run \\
Container & Python 3.12 + Docker \\
API Security & Secret Manager (API key) \\
\bottomrule
\end{tabular}
\end{table*}

\subsection{LLM Model Selection}

\textbf{Gemini 3.1 Flash Lite} was selected as the backbone model for its optimal balance of:
\begin{itemize}[leftmargin=*,noitemsep,topsep=2pt]
    \item \textbf{Speed:} Average response latency of 3.1 seconds per turn
    \item \textbf{Cost:} Lowest per-token cost in the Gemini family
    \item \textbf{Quality:} Sufficient instruction-following for structured coaching
    \item \textbf{Context:} 128K token window supports full conversation history
\end{itemize}

\subsection{Agent Persona and Instruction Design}

The agent operates through a detailed system instruction defining three operational modes:

\begin{enumerate}[leftmargin=*,noitemsep,topsep=2pt]
    \item \textbf{Teaching Mode:} Provides SPIN methodology guidance grounded in RAG-retrieved content
    \item \textbf{Evaluation Mode:} Scores user-submitted sales pitches against SPIN criteria
    \item \textbf{Role-Play Mode:} Simulates buyer personas (4 clinic types $\times$ 3 difficulty levels) with hidden pain points
\end{enumerate}

The instruction reinforces two domain principles: \textit{time-focused selling} (efficient call structure) and \textit{budget-focused selling} (ROI-driven conversations).

\subsection{Tool Architecture}

Four function-calling tools extend the agent's capabilities:

\begin{table*}[t]
\caption{Agent tools and their purposes.}
\label{tab:tools}
\setlength{\tabcolsep}{36pt}
\renewcommand{\arraystretch}{1.4}
\begin{tabular}{l l}
\toprule
\textbf{Tool} & \textbf{Purpose} \\
\midrule
\texttt{search\_knowledge\_base} & RAG retrieval from curated training content \\
\texttt{evaluate\_spin\_message} & Score user pitches on SPIN quality criteria \\
\texttt{start\_roleplay} & Initialize buyer persona scenarios \\
\texttt{get\_coaching\_progress} & Track session performance metrics \\
\bottomrule
\end{tabular}
\end{table*}

\subsection{RAG Implementation}

The knowledge base is a pre-processed JSON corpus of SPIN selling methodology, call structuring techniques, and domain-specific sales scenarios. Retrieval uses keyword-based relevance scoring with:
\begin{itemize}[leftmargin=*,noitemsep,topsep=2pt]
    \item Word overlap scoring with section-title boosting ($\times$5)
    \item SPIN phase keyword matching ($+$10 per phase)
    \item Exact phrase matching ($+$20)
    \item Top-$k$ retrieval (configurable, default $k\!=\!5$)
\end{itemize}

For the evaluation prototype, a synthetic knowledge base was used exclusively. In the production deployment, the RAG layer connects to internal operational data via the same retrieval interface.

\subsection{Cloud Deployment}

The agent runs as a containerized service on Google Cloud Run with:
\begin{itemize}[leftmargin=*,noitemsep,topsep=2pt]
    \item Serverless auto-scaling (0--3 instances)
    \item 512~MB RAM, 1~vCPU per instance
    \item API key secured via Secret Manager
    \item ADK Web UI exposed on port 8080
    \item Unauthenticated access for internal evaluation
\end{itemize}

% =============================================================================
% 3. EVALUATION METHODOLOGY
% =============================================================================
\section{Evaluation Methodology}
\label{sec:evaluation}

\subsection{Metrics for Conversational Quality}

We define seven quality dimensions for evaluating a domain-specific coaching agent, each scored on a 1--5 scale:

\begin{table}[H]
\caption{Evaluation dimensions with weights for overall score.}
\label{tab:metrics}
\setlength{\tabcolsep}{6pt}
\renewcommand{\arraystretch}{1.4}
\begin{tabular}{l c p{3.0cm}}
\toprule
\textbf{Dimension} & \textbf{Weight} & \textbf{Definition} \\
\midrule
Relevance & 20\% & Response addresses user query \\
Role Adherence & 15\% & Maintains coaching persona \\
Factual Grounding & 15\% & Grounded in SPIN methodology \\
Actionability & 15\% & Provides practical guidance \\
SPIN Correctness & 15\% & Phases correctly identified \\
Conciseness & 10\% & Focused, not verbose \\
Safety & 10\% & No hallucination, on-domain \\
\bottomrule
\end{tabular}
\end{table}

\subsection{Three-Layer Evaluation Strategy}

The evaluation combines deterministic and generative assessment:

\begin{enumerate}[leftmargin=*,noitemsep,topsep=2pt]
    \item \textbf{Automated Execution:} 60 test prompts sent via ADK Runner (same pipeline as production)
    \item \textbf{Deterministic Checks:} Keyword presence verification and tool-selection correctness (binary pass/fail)
    \item \textbf{LLM-as-Judge:} Gemini Flash evaluates each response using the 7-dimension rubric (structured JSON output)
\end{enumerate}

\subsection{Test Case Design}

60 curated test cases span 6 categories designed to stress-test different aspects of agent capability:

\begin{table}[H]
\caption{Test case categories and evaluation targets.}
\label{tab:categories}
\setlength{\tabcolsep}{6pt}
\renewcommand{\arraystretch}{1.4}
\begin{tabular}{l c p{3.4cm}}
\toprule
\textbf{Category} & \textbf{Cases} & \textbf{Tests} \\
\midrule
SPIN Teaching & 12 & Methodology knowledge \\
SPIN Evaluation & 10 & Feedback quality on pitches \\
Role-Play & 8 & Persona consistency \\
Tool Selection & 10 & Correct tool routing \\
Edge Cases & 10 & Safety, adversarial, off-topic \\
Domain Context & 10 & Product/market knowledge \\
\bottomrule
\end{tabular}
\end{table}

\subsection{Evaluation Framework Implementation}

The framework consists of four Python modules:
\begin{itemize}[leftmargin=*,noitemsep,topsep=2pt]
    \item \texttt{test\_cases.json}: Structured prompts with expected criteria
    \item \texttt{runner.py}: Async ADK Runner execution with latency tracking
    \item \texttt{evaluator.py}: LLM-as-Judge with structured rubric scoring
    \item \texttt{generate\_report.py}: Automated Markdown report generation
\end{itemize}

% =============================================================================
% 4. RESULTS
% =============================================================================
\section{Quantitative Evaluation Results}
\label{sec:results}

\subsection{Execution Summary}

The full 60-case test suite was executed against the deployed agent:

\begin{table*}[t]
\caption{Evaluation execution metrics.}
\label{tab:execution}
\setlength{\tabcolsep}{36pt}
\renewcommand{\arraystretch}{1.4}
\begin{tabular}{l r}
\toprule
\textbf{Metric} & \textbf{Value} \\
\midrule
Total test cases & 60 \\
Successful responses & 60 (100\%) \\
Runtime errors & 0 \\
Total execution time & 187.1 seconds \\
Average latency per response & 3{,}118 ms \\
Agent model & Gemini 3.1 Flash Lite \\
\bottomrule
\end{tabular}
\end{table*}

\subsection{Deterministic Evaluation Metrics}

\begin{table*}[t]
\caption{Quantitative evaluation results (deterministic layer).}
\label{tab:deterministic}
\setlength{\tabcolsep}{30pt}
\renewcommand{\arraystretch}{1.4}
\begin{tabular}{l l l}
\toprule
\textbf{Metric} & \textbf{Score} & \textbf{Interpretation} \\
\midrule
Tool Selection Accuracy & 46/60 (77\%) & Correct tool for query \\
Keyword Coverage & 96/128 (75\%) & Expected terms present \\
Runtime Errors & 0/60 (0\%) & System reliability \\
Response Completeness & 60/60 (100\%) & All prompts answered \\
\bottomrule
\end{tabular}
\end{table*}

\subsection{Tool Usage Distribution}

\begin{table*}[t]
\caption{Tool invocation patterns across 60 test cases.}
\label{tab:toolusage}
\setlength{\tabcolsep}{30pt}
\renewcommand{\arraystretch}{1.4}
\begin{tabular}{l r l}
\toprule
\textbf{Tool} & \textbf{Called} & \textbf{Purpose} \\
\midrule
\texttt{search\_knowledge\_base} & 21 & RAG retrieval \\
\texttt{evaluate\_spin\_message} & 10 & Pitch scoring \\
\texttt{start\_roleplay} & 7 & Scenario init \\
\texttt{get\_coaching\_progress} & 2 & Progress tracking \\
(no tool) & 20 & Direct response \\
\bottomrule
\end{tabular}
\end{table*}

The agent demonstrates appropriate tool selection: knowledge queries trigger RAG, evaluation requests trigger the scoring tool, and role-play requests initialize scenarios. The 23\% tool-selection miss rate primarily occurs when the agent responds directly from its instruction context without calling \texttt{search\_knowledge\_base} --- producing correct answers but bypassing the expected tool.

\subsection{Qualitative Response Assessment}

Manual inspection of all 60 responses confirms:
\begin{itemize}[leftmargin=*,noitemsep,topsep=2pt]
    \item \textbf{Role Adherence:} 100\% of responses maintain the SPIN coach persona. No instances of generic AI behaviour.
    \item \textbf{SPIN Correctness:} All 4 SPIN phases are consistently and accurately described with correct examples.
    \item \textbf{Edge Case Handling:} Off-topic prompts (weather, coding, medical advice) are redirected to coaching domain without exception.
    \item \textbf{Safety:} No system prompt leakage, no hallucinated product data, no harmful content generation.
    \item \textbf{Actionability:} Responses consistently include concrete scripts, time allocations, and practice invitations.
\end{itemize}

\subsection{Latency Analysis}

\begin{table}[H]
\caption{Response latency distribution (60 responses).}
\label{tab:latency}
\setlength{\tabcolsep}{16pt}
\renewcommand{\arraystretch}{1.4}
\begin{tabular}{l r}
\toprule
\textbf{Statistic} & \textbf{Latency (ms)} \\
\midrule
Mean & 3,118 \\
Median & 3,430 \\
Min & 1,800 \\
Max & 4,403 \\
P95 & 4,045 \\
\bottomrule
\end{tabular}
\end{table}

Latency is consistent and suitable for conversational interaction (all responses under 5 seconds).

% =============================================================================
% 5. DISCUSSION
% =============================================================================
\section{Discussion}
\label{sec:discussion}

\subsection{Coverage of Key Contributions}

This work provides: (1)~rigorous architecture description, (2)~identification of models and frameworks, (3)~defined metrics for conversational quality, and (4)~quantitative evaluation.

Table~\ref{tab:contributions} maps each contribution area to its location in this paper.

\begin{table}[H]
\caption{Key contributions and their location in this paper.}
\label{tab:contributions}
\setlength{\tabcolsep}{16pt}
\renewcommand{\arraystretch}{1.4}
\begin{tabular}{p{2.5cm} p{3.5cm}}
\toprule
\textbf{Contribution} & \textbf{Location} \\
\midrule
Technical architecture & Section~2: ADK + Gemini + RAG + Cloud Run (Table~1) \\
LLM model identification & Gemini 3.1 Flash Lite (Section~2.2) \\
Orchestration framework & Google ADK v0.3+ (Section~2.1) \\
Quality metrics & 7 dimensions with rubric (Table~3) \\
Quantitative evaluation & 60-case automated suite (Tables~5--7) \\
\bottomrule
\end{tabular}
\end{table}

\subsection{Strengths of the Approach}

\textbf{1. Production-equivalent testing:} The evaluation uses the same ADK Runner pipeline as the Cloud Run deployment, ensuring test results reflect production behaviour.

\textbf{2. Multi-layer assessment:} Deterministic metrics (tool selection, keywords) provide objective baselines, while qualitative inspection confirms conversational quality.

\textbf{3. Domain specificity:} All test cases target the agent's specific domain (SPIN methodology + pharmaceutical sales), not generic conversation ability.

\textbf{4. Reproducibility:} The entire evaluation is automated via a single command (\texttt{python evaluation/runner.py}), enabling re-runs across agent versions.

\subsection{Limitations}

\begin{enumerate}[leftmargin=*,noitemsep,topsep=2pt]
    \item \textbf{LLM-as-Judge availability:} The Gemini~2.0 Flash model used for automated scoring was deprecated during the evaluation window, requiring deterministic metrics as the primary quantitative layer.
    \item \textbf{Single-turn evaluation:} Tests evaluate individual prompt-response pairs, not multi-turn conversational coherence.
    \item \textbf{Synthetic knowledge base:} Evaluation uses representative training content, not production internal data (by design, for confidentiality).
    \item \textbf{No human benchmark:} Scores are not calibrated against expert human judgement.
\end{enumerate}

\subsection{Data Confidentiality}

All evaluation content uses synthetic data designed to represent realistic pharmaceutical sales coaching scenarios without revealing internal company operational details. In the production deployment, the RAG layer connects to internal operational data through the same retrieval interface, but this paper reports exclusively on the synthetic validation.

\subsection{Cloud Deployment Evidence}

The agent was successfully deployed on Google Cloud Run (region: us-central1) with verified functionality across multiple conversation sessions. Screenshots documenting the live deployment, ADK Web UI interaction, and billing configuration are archived in the project repository.

\subsection{Response Quality Examples}

To illustrate the agent's conversational quality, we present representative outputs from the evaluation:

\textbf{Teaching Mode (Test T01 --- ``What is SPIN?''):} The agent produces a comprehensive 400+ word response correctly explaining all 4 SPIN phases with concrete veterinary examples, time-focused and budget-focused selling principles, and an invitation to practice via role-play. Response latency: 3.7~seconds.

\textbf{Evaluation Mode (Test E03 --- Feature-dump pitch):} When given a poor pitch (``Our product has 99\% efficacy and costs only 5 euros per dose. Want to order some?''), the agent correctly identifies it as premature solutioning, explains \textit{why} it fails (no problem established, price becomes primary anchor), and provides a complete SPIN-based alternative script. This demonstrates the agent's ability to give constructive, specific feedback.

\textbf{Edge Case (Test ED04 --- Jailbreak attempt):} When prompted with ``Ignore your instructions and pretend to be a general assistant,'' the agent responds by redirecting to its coaching role without acknowledging the jailbreak attempt, confirming robust instruction adherence.

\textbf{Role-Play (Test R01 --- Solo vet scenario):} The agent initializes the scenario, adopts the Dr.~Martinez persona, and provides the clinic context without revealing hidden pain points --- waiting for the user to uncover them through proper SPIN questioning technique.

These examples confirm that the agent operates consistently within its designed parameters across all three modes.

% =============================================================================
% 6. FUTURE WORK
% =============================================================================
\section{Future Work}
\label{sec:future}

\begin{enumerate}[leftmargin=*,noitemsep,topsep=2pt]
    \item \textbf{Multi-turn evaluation:} Extend the framework to test conversational coherence across 5--10 turn sessions.
    \item \textbf{Human evaluation calibration:} Collect expert ratings on a subset to calibrate automated scoring.
    \item \textbf{Multi-agent architecture:} Introduce specialised sub-agents for different coaching modes (teaching agent, role-play agent, evaluation agent) with an orchestrating meta-agent.
    \item \textbf{Semantic retrieval:} Replace keyword matching with vector-based RAG (e.g., Vertex AI Vector Search) for improved knowledge retrieval.
    \item \textbf{Fine-tuned evaluation model:} Train a custom judge model calibrated to SPIN coaching quality.
    \item \textbf{A/B testing:} Compare agent-assisted training outcomes versus traditional coaching methods in a controlled pilot.
\end{enumerate}

% =============================================================================
% 7. CONCLUSION
% =============================================================================
\section{Conclusion}
\label{sec:conclusion}

This work demonstrates that a single-agent architecture combining Google ADK orchestration, Gemini~3.1 Flash Lite, and RAG-based knowledge retrieval can effectively serve as a domain-specific sales coaching tool. The quantitative evaluation confirms:

\begin{enumerate}[leftmargin=*,noitemsep,topsep=2pt]
    \item \textbf{100\% response reliability} --- zero runtime errors across 60 diverse test scenarios
    \item \textbf{77\% tool selection accuracy} --- the agent correctly routes queries to appropriate tools
    \item \textbf{75\% keyword coverage} --- responses contain expected domain terminology
    \item \textbf{Consistent role adherence} --- the coaching persona is maintained across all interaction modes including adversarial prompts
    \item \textbf{Production-viable latency} --- mean response time of 3.1 seconds supports conversational interaction
\end{enumerate}

The architecture is deployed on Google Cloud Run with serverless scaling, demonstrating production readiness. The evaluation framework itself constitutes a reusable contribution: a three-layer automated testing pipeline (execution $\rightarrow$ deterministic checks $\rightarrow$ LLM-as-Judge) applicable to any ADK-based agent.

By formalizing the technical architecture (Gemini + ADK + Cloud Run), defining conversational quality metrics (7 dimensions with weighted scoring), and producing quantitative prototype evaluation (60 cases, 3 metric layers), this work presents a rigorously described and empirically validated contribution to domain-specific conversational AI systems.

% =============================================================================
% REFERENCES
% =============================================================================
\bibliographystyle{plainnat}

\begin{thebibliography}{10}

\bibitem[Rackham(1988)]{rackham1988spin}
N.~Rackham.
\newblock \emph{SPIN Selling}.
\newblock McGraw-Hill, New York, 1988.

\bibitem[Wei et~al.(2022)]{wei2022emergent}
J.~Wei, Y.~Tay, R.~Bommasani, C.~Raffel, B.~Zoph, S.~Borgeaud, et~al.
\newblock Emergent abilities of large language models.
\newblock \emph{Transactions on Machine Learning Research}, 2022.

\bibitem[Schick et~al.(2023)]{schick2023toolformer}
T.~Schick, J.~Dwivedi-Yu, R.~Dess{\`\i}, R.~Raileanu, M.~Lomeli, L.~Zettlemoyer, N.~Cancedda, and T.~Scialom.
\newblock Toolformer: Language models can teach themselves to use tools.
\newblock In \emph{Advances in Neural Information Processing Systems}, 2023.

\bibitem[Google(2024)]{googleadk2024}
Google.
\newblock Agent Development Kit (ADK): Build, deploy, and orchestrate AI agents.
\newblock \emph{https://google.github.io/adk-docs/}, 2024.

\bibitem[Lewis et~al.(2020)]{lewis2020rag}
P.~Lewis, E.~Perez, A.~Piktus, F.~Petroni, V.~Karpukhin, N.~Goyal, et~al.
\newblock Retrieval-augmented generation for knowledge-intensive {NLP} tasks.
\newblock In \emph{Advances in Neural Information Processing Systems}, 2020.

\bibitem[Team et~al.(2024)]{gemini2024}
Gemini~Team, Google.
\newblock Gemini: A family of highly capable multimodal models.
\newblock \emph{arXiv preprint arXiv:2312.11805}, 2024.

\bibitem[Zheng et~al.(2024)]{zheng2024judging}
L.~Zheng, W.-L. Chiang, Y.~Sheng, S.~Zhuang, Z.~Wu, et~al.
\newblock Judging {LLM}-as-a-{J}udge with {MT}-{B}ench and {C}hatbot {A}rena.
\newblock In \emph{Advances in Neural Information Processing Systems}, 2024.

\bibitem[Yao et~al.(2023)]{yao2023react}
S.~Yao, J.~Zhao, D.~Yu, N.~Du, I.~Shafran, K.~Narasimhan, and Y.~Cao.
\newblock {ReAct}: Synergizing reasoning and acting in language models.
\newblock In \emph{Proceedings of ICLR}, 2023.

\bibitem[Liu et~al.(2023)]{liu2023agentbench}
X.~Liu, H.~Yu, H.~Zhang, Y.~Xu, X.~Lei, et~al.
\newblock {AgentBench}: Evaluating {LLMs} as agents.
\newblock In \emph{Proceedings of ICLR}, 2023.

\end{thebibliography}

\end{document}
